\lecture{5}{Tue 16 Sep 2025 12:32}{Covariance}

Before discussing covariance, we continue our discussion of the Equivalence principle. Consider a laser firing a photon at a fixed wavelength $\lambda _i$. We measure it at the top with a detector to be a \emph{different number} $\lambda _f \neq \lambda _i$. This experiment was done at Harvard by shooting photons 22 meters. This small difference in gravity was able to measure the redshift. We derive this using only the equivalence principle.

This situation is equivalent to the entire tower accelerating upwards at $9.8m/s^2$, and shooting a laser up the tower the instant it starts accelerating. The distance $z(t)$ of the photon from $z_0$ is $z(t) = z_0+\frac{1}{2} at^2$. Of course, $\Delta t=t_f-t_i=\frac{z_f}{c} $. We choose $t_i=0$. The velocity of the rocket is simply $v=at_f=a\frac{z_0}{c} $. There is then a redshift factor of $\frac{\Delta \lambda }{\lambda _i}  $ by the Doppler effect. This gives a doppler shift of
\[
	\sqrt{\frac{1+v/c}{1-v/c} } -1
.\]
This approximates out to $az_f/c^2$ in the small gravitational potential when plugging in the $t$s.

We use a mathematical systemization of this argument, called \emph{general covariance}.

\begin{enumerate}
	\item First, write down equations of motion for physics in the absense of gravity.
	\item Then, introduce the metric and its derivatives into the equations of motion in such a way that the resultant equations of motion are invariant under coordinate changes (Lorentz transformations). We then WLOG reduce to Minkowski space.
\end{enumerate}

To do this, we need to learn how to efficiently write the equations of motion in a coordinate-invariant manner more efficiently. We will explore this at length. The fascinating implications of this fact are that gravity follows completely by coordinate-invariance and the Minkowski metric, meaning the \emph{entire effect of gravity} is the coordinate transformation into the gravitational frame.

We start by considering free motion of a particle
\[
	\frac{\mathrm{d}u^\mu }{\mathrm{d}\tau } =0
.\]
Recall that $u^\mu $ is  the 4-velocity $(u^0,u^i)$ and $\tau $ is the proper time (rest frame time)
\[
	\mathrm{d} \tau ^2=-\eta _{\mu \nu } \mathrm{d} x^\mu \mathrm{d} x^\nu 
.\]
. Idk $\tau $ represents time in our coordinate system. This is equivalently written as
\[
	\frac{\mathrm{d}^2x^\mu }{\mathrm{d}\tau ^2} =\left( \frac{\mathrm{d}^2t}{\mathrm{d}\tau ^2} ,\frac{\mathrm{d}^2x^i}{\mathrm{d}\tau ^2}  \right) =0
.\]
In the rest frame, $t=\tau $, so the time component vanishes in the rest frame.

The issue is we want to also work with massless particles, and they don't have proper time. This means we can't use $\tau $ since it never changes, so we have to do something else. Define $\sigma =x^0$ as the time component in an inertial frame that is not necessarily the reference frame. We then have
\[
	\frac{\mathrm{d}^2x^\mu }{\mathrm{d}\sigma ^2} =0
.\]
We use affine parameters to show that $\sigma $ is invariant under our choice. Write $\sigma =a\lambda +b$, where $\lambda $ is an affine parameter. We then still have
\[
	\frac{\mathrm{d}^2x^\mu }{\mathrm{d}\lambda ^2} =0
\]
by change of variables. We thus have the freedom to redefine the coordinate system using any affine parameter.

Both massive \emph{and} massless particles can then be written as
\[
	\frac{\mathrm{d}^2\xi ^\mu }{\mathrm{d}\lambda ^2} =0,\qquad \xi ^\mu =\text{freely falling frame}.
\]
Of course by Lorentz invariance (and in particular, invariance under boosts), we see that the velocity of the particle in this frame need not be defined, since all are Lorentz invariant.

The question is: what are the equations of motion in the lab frame? The answer is given by transforming
\[
	\frac{\mathrm{d}^2\xi ^\mu }{\mathrm{d}\lambda ^2} =0 \to \text{lab frame equivalent} 
.\]
We do this explicitly by computing $\xi ^\mu (x)$. We want to write these in terms of the metric. In the freely falling frame we choose the metric by $\eta _{\mu \nu } $ with vanishing first derivatives. This is what the freely falling frame is (recall it exists by earlier discussion). The lab frame is an accelerating frame because our lab is on a rocketship. I love GR.

Write $g_{\mu \nu } $ for the metric for the lab frame, and recall that we have that our metric transformation is given by
\[
	g_{\mu \nu } \to g'_{\mu \nu } =g_{\alpha \beta } \frac{\partial x^\alpha }{\partial (x')^\mu } \frac{\partial x^\beta }{\partial (x')^\nu } =\partial _\mu x^\alpha \partial _\nu x'^\beta =g_{\alpha \beta } \frac{\partial x^\alpha }{\partial \xi ^\mu } \frac{\partial x^\beta }{\partial \xi ^\nu } 
.\]
A trick to remembering this is that the indices must always contract meaning since $\mu ,\nu $ are on the bottom for $g_{\mu \nu } $, the must also be on the bottom, which we have by $\partial _\mu ,\partial _\nu $.

Since $g'$ is the freely falling frame, we have
\[
	g_{\alpha \beta } \frac{\partial x^\alpha }{\partial \xi ^\mu } \frac{\partial x^\beta }{\partial \xi ^\nu } =\eta _{\mu \nu } 
.\]
Writing the tensors as matrices, we have inverse matrices
\[
	\frac{\partial x^\alpha }{\partial \xi ^\mu } =\left( \frac{\partial \xi ^\alpha }{\partial x^\mu }  \right) ^{-1} 
.\]
Multiplying through on our above equation gives
\[
	\eta \frac{\partial \xi ^\mu }{\partial x^\sigma } =g_{\alpha \beta } \frac{\partial x^\alpha }{\partial \xi ^\mu } \frac{\partial \xi ^\mu }{\partial x^\sigma } \frac{\partial x^\beta }{\partial \xi ^\nu } =g_{\alpha \beta } \delta _\sigma ^\alpha \frac{\partial x^\beta }{\partial \xi ^\nu } =g_{\sigma \beta } \frac{\partial x^\beta }{\partial x^\nu } 
.\]
Applying the same argument again gives
\[
	\eta _{\mu \nu } \frac{\partial \xi ^\mu }{\partial x^\sigma } \frac{\partial \xi ^\nu }{\partial x^\rho } =g_{\sigma \rho } 
.\]

We are now writing $\tau $ instead of $\lambda $ because the prof's notation got lazy. Applying the chain rule and product rule,
\[
	0=\frac{\mathrm{d}^2\xi ^\mu }{\mathrm{d}\tau ^2} =\frac{\mathrm{d}}{\mathrm{d}\tau } \frac{\mathrm{d}}{\mathrm{d}\tau } \xi ^\mu =\frac{\mathrm{d}}{\mathrm{d}\tau } \left( \frac{\mathrm{d}x^\alpha }{\mathrm{d}\tau } \frac{\partial \xi ^\mu }{\partial x^\alpha }  \right) =\frac{\partial ^2\xi ^\mu }{\partial x^\alpha \partial x^\beta } \frac{\mathrm{d}x^\alpha }{\mathrm{d}\tau } \frac{\mathrm{d}x^\beta }{\mathrm{d}\tau } +\frac{\mathrm{d}^2x^\alpha }{\mathrm{d}\tau ^2} \frac{\partial \xi ^\mu }{\partial x^\alpha } 
.\]
Of course throughout this we are assuming derivatives freely commute, which indeed follows for smooth maps. We use the same process to strip of the factor of $\partial_\alpha \xi ^\mu $ by multiplying through by $\partial_{\xi ^\mu } x^\sigma $:
\[
	0=\frac{\mathrm{d}x^\sigma }{\mathrm{d}\tau ^2} +\frac{\mathrm{d}x^\alpha }{\mathrm{d}\tau } \frac{\mathrm{d}x^\beta }{\mathrm{d}\tau } \frac{\partial x^\sigma }{\partial \xi ^\mu }\frac{\partial ^2\xi ^\mu }{\partial x^\alpha \partial x^\beta } 
.\]
Wr write
\[
	\Gamma _{\alpha \beta } ^\sigma  \coloneqq \frac{\partial ^2\xi ^\mu }{\partial x^\alpha \partial x^\beta } \frac{\partial x^\sigma }{\partial \xi ^\mu } 
.\]
The answer is then rewritten as
\[
	0=\frac{\mathrm{d}^2x^\sigma }{\mathrm{d}\tau ^2} +\Gamma _{\alpha \beta } ^\sigma \frac{\mathrm{d}x^\alpha }{\mathrm{d}\tau } \frac{\mathrm{d}x^\beta }{\mathrm{d}\tau } 
.\]
This is called the \emph{geodesic equation}. We want to rewrite it in terms of a metric to be able to use it more easily. To do this, we have to rewrite $\Gamma _{\alpha \beta } ^\sigma $ in terms of the metric
\[
	\eta _{\mu \nu } \frac{\partial \xi ^\mu }{\partial x^\sigma } \frac{\partial \xi ^\nu }{\partial x^\rho } =g_{\sigma \rho } 
.\]
Differentiate with respect to $x^\lambda $:
\[
	\frac{\partial g_{\alpha \beta } }{\partial x^\lambda } =\eta _{\mu \nu } \frac{\partial ^2\xi ^\mu }{\partial x^\lambda \partial x^\alpha } \frac{\partial \xi ^\nu }{\partial x^\beta } +\eta _{\mu \nu } \frac{\partial \xi ^\mu }{\partial x^\alpha } \frac{\partial ^2\xi ^\nu }{\partial x^\lambda \partial x^\beta } 
.\]
Notice that this is simply
\[
	\frac{\partial g_{\mu \nu } }{\partial x^\lambda } =\eta _{\mu \nu } \left( \Gamma _{\lambda \alpha } ^\rho \frac{\partial \xi ^\mu }{\partial x^\rho } \frac{\partial \xi ^\nu }{\partial x^\beta } +\Gamma _{\lambda \beta } ^\rho \frac{\partial \xi ^\mu }{\partial x^\alpha } \frac{\partial \xi ^\nu }{\partial x^\rho }  \right) =\Gamma^\rho_{\lambda \alpha } g_{\rho \beta } +\Gamma ^\rho _{\lambda \beta } g_{\alpha \rho } ,
\]
where the final step follows also by our above formula for $g_{\mu \nu } $.

Note that by definition, the Christofell symbol $\Gamma _{\alpha \beta } ^\mu $ is symmetric under the lower indices, so we add the equation to itself under a permutation of indices, which I don't want tow rite. We end up with
\[
	\partial _\lambda g_{\mu \nu } +\partial _\mu g_{\lambda \nu } -\partial _\nu g_{\mu \lambda } =2\Gamma ^\rho _{\lambda \mu } g_{\rho \nu }  
.\]
Stripping the metric by multiplying by its inverse:
\[
	\Gamma _{\alpha \beta } ^\rho =\frac{1}{2} g^{\sigma \rho } \left( \frac{\partial g_{\alpha \sigma } }{\partial x^\beta } +\frac{\partial g_{\beta \sigma } }{\partial x^\alpha } -\frac{\partial g_{\alpha \beta } }{\partial x^\sigma }  \right) 
.\]
This allows us to rewrite the geodesic equation in terms of the metric only. Since the Minkowski metric has fully dropped out, this olds in \emph{any metric}.
